\chapter{Experiment Configuration Catalogue}
\label{app:config_catalogue}

Not all configurations listed here are discussed equally in the main text;
the thesis chapters focus on a representative subset of nine benchmark
regimes chosen for interpretability and computational completeness.
The parameter values for each dynamics regime are given in
Tables~\ref{tab:ndyn_regimes_full} and~\ref{tab:do_regimes_full}
(Appendix~\ref{app:full_regime_catalogue}).  This appendix lists every
YAML file in \texttt{configs/systematic/} for reference.

\noindent\footnotesize
Most NDYN base files have a \emph{\_VS} companion (\cv{speed\_mode}${}=$\texttt{variable}), except where the speed mode is already variable or no companion was defined; see Appendix~\ref{app:full_regime_catalogue} for exceptions.
All DO configs share $v_0{=}3$, $\eta{=}0.2$, $C_a{=}1$, $\ell_a{=}1$; only $C_r,\ell_r$ vary.

\medskip
\begin{minipage}[t]{0.48\textwidth}\raggedright
\captionof{table}{NDYN configs (physics-knob sweep).}
\renewcommand{\arraystretch}{0.92}
\begin{tabular}{@{}ll@{}}
\toprule
File & Regime \\
\midrule
\cv{NDYN01\_crawl}           & Slow, low noise \\
\cv{NDYN02\_flock}           & Coherent flock \\
\cv{NDYN03\_sprint}          & High speed \\
\cv{NDYN04\_gas}             & High noise \\
\cv{NDYN05\_blackhole}       & Strong attraction \\
\cv{NDYN06\_supernova}       & Dominant repulsion \\
\cv{NDYN07\_crystal}         & Balanced forces \\
\cv{NDYN08\_pure\_vicsek}    & No Morse forces \\
\cv{NDYN09\_longrange}       & Large length scales \\
\cv{NDYN10\_shortrange}      & Small length scales \\
\cv{NDYN11\_noisy\_collapse} & Noise + attraction \\
\cv{NDYN12\_fast\_explosion} & Speed + repulsion \\
\cv{NDYN13\_chaos}           & Extreme speed/noise \\
\cv{NDYN14\_varspeed}        & Variable speed \\
\bottomrule
\end{tabular}
\end{minipage}\hfill
\begin{minipage}[t]{0.48\textwidth}\raggedright
\captionof{table}{DO configs (outcome-targeted sweep).}
\renewcommand{\arraystretch}{0.92}
\begin{tabular}{@{}ll@{}}
\toprule
File & Target pattern \\
\midrule
\cv{DO\_CS01\_swarm\_C01\_l05}     & Cohesive swarm \\
\cv{DO\_CS02\_swarm\_C05\_l3}      & Cohesive swarm \\
\cv{DO\_CS03\_swarm\_C09\_l3}      & Cohesive swarm \\
\cv{DO\_DM01\_dmill\_C09\_l05}     & Double mill \\
\cv{DO\_DR01\_dring\_C01\_l01}     & Double ring \\
\cv{DO\_DR02\_dring\_C09\_l09}     & Double ring \\
\cv{DO\_EC01\_esccol\_C2\_l3}      & Escape--collapse \\
\cv{DO\_EC02\_esccol\_C3\_l05}     & Escape--collapse \\
\cv{DO\_ES01\_escsym\_C3\_l09}     & Escape symmetric \\
\cv{DO\_EU01\_escuns\_C2\_l2}      & Escape unstable \\
\cv{DO\_EU02\_escuns\_C3\_l3}      & Escape unstable \\
\cv{DO\_SM01\_mill\_C05\_l01}      & Single mill \\
\cv{DO\_SM02\_mill\_C3\_l01}       & Single mill \\
\cv{DO\_SM03\_mill\_C2\_l05}       & Single mill \\
\bottomrule
\end{tabular}
\end{minipage}\normalsize

\noindent Including the \_VS variants, the repository contains
\textbf{54 YAML files} in \texttt{configs/systematic/}.
Six of these (three near-duplicate base regimes, NDYN03, NDYN11,
NDYN12, and their VS counterparts) were removed before running,
leaving 48 candidate configurations; of those, \textbf{37} produced
usable results and form the catalogue analysed in Chapters~7--9
(Appendix~\ref{app:full_regime_catalogue}).
The repository contains the full configuration set used during
development, including duplicate or superseded files.  The systematic
catalogue analysed in this thesis retains 37 regimes after removing
duplicate NDYN entries; of these, nine form the primary benchmark used
in the main-text comparisons, while the remainder are used for
catalogue-level robustness checks.
Each configuration is run across \textbf{4 IC families} $\times$
multiple runs per family, totalling several thousand independent
simulations.

% -----------------------------------------------------------------------
\section{Configuration Schema}
\label{app:config_schema}
% -----------------------------------------------------------------------

Every experiment is fully specified by a single YAML file read at boot time.
The schema below mixes global defaults with regime-specific
overrides; only a subset of keys varies across the systematic regime files
listed above.
The schema mirrors an actual experiment configuration,
annotated with allowed values.
The EF-ROM keys (\cv{sim} through \cv{eval}) follow
\cv{NDYN02\_flock.yaml}; the \cv{wsindy} block follows
\cv{DYN1\_gentle\_wsindy.yaml}.

Each YAML file contains twelve top-level sections:
\texttt{sim} (agent count, training horizon, domain geometry),
\texttt{test\_sim} (forecast horizon),
\texttt{model} (agent dynamics type and speed mode),
\texttt{noise} (angular noise magnitude),
\texttt{forces} (Morse attraction/repulsion parameters $C_a$, $C_r$, $\ell_a$, $\ell_r$),
\texttt{alignment} (Vicsek radius and coupling strength),
\texttt{train\_ic} and \texttt{test\_ic} (initial-condition family and distributional parameters),
\texttt{rom} (POD rank, MVAR lag, LSTM architecture and training settings),
\texttt{eval} (forecast metrics, bootstrap replicates, confidence level),
\texttt{wsindy} (library scope, thresholding schedule, test-function parameters),
and \texttt{experiment\_name} (unique run identifier).
Full annotated schemas for all nine benchmark regimes are
available in the repository at
\url{https://github.com/eugenio-macias/wsindy-manifold/tree/main/configs}.

% -----------------------------------------------------------------------
\section{Symbol-to-Configuration Key Mapping}
\label{app:symbol_config_map}
% -----------------------------------------------------------------------

Table~\ref{tab:symbol_config} maps the mathematical symbols used in the
main text to their corresponding YAML configuration keys and source
modules in the accompanying repository.

\begin{table}[htbp]
\centering
\caption{Mapping from mathematical symbols to configuration keys
and source modules used in the code repository.}
\label{tab:symbol_config}
\renewcommand{\arraystretch}{1.05}
\small
\begin{tabular}{@{}lll@{}}
\toprule
\textbf{Math symbol / quantity}  & \textbf{Config key} & \textbf{Source module} \\
\midrule
$\Delta t$ (time step)                & \cv{sim.dt}                     & \cv{sim\_runner.py}            \\
$N$ (agent count)                     & \cv{sim.traj}                   & \cv{sim\_runner.py}            \\
$v_0$ (speed)                         & \cv{model.speed}                & \cv{particle\_model.py}        \\
$R$ (alignment radius)                & \cv{alignment.R}                & \cv{particle\_model.py}        \\
$\eta$ (noise strength)               & \cv{noise.eta}                  & \cv{particle\_model.py}        \\
Noise model (uniform/wrapped-Cauchy)  & \cv{noise.kind}                 & \cv{particle\_model.py}        \\
Match variance flag                   & \cv{noise.match\_variance}      & \cv{particle\_model.py}        \\
KDE bandwidth                         & \cv{rom.density\_bandwidth}     & \cv{pod\_builder.py}           \\
Shift-alignment reference             & \cv{rom.shift\_align\_ref}      & \cv{shift\_align.py}           \\
Density transform ($\sqrt{\cdot}$)    & \cv{rom.density\_transform}     & \cv{pod\_builder.py}           \\
$r$ (POD rank)                        & \cv{rom.fixed\_modes}           & \cv{pod\_builder.py}           \\
Clamp mode (C1)                       & \cv{rom.clamp\_mode}            & \cv{test\_evaluator.py}        \\
Mass post-processing                  & \cv{rom.mass\_postprocess}      & \cv{test\_evaluator.py}        \\
MVAR lag $p$                          & \cv{rom.mvar.lag}               & \cv{mvar\_trainer.py}          \\
Ridge penalty $\alpha$                & \cv{rom.mvar.ridge\_alpha}      & \cv{mvar\_trainer.py}          \\
Eigenvalue threshold                  & \cv{rom.eigenvalue\_threshold}  & \cv{mvar\_trainer.py}          \\
LSTM subsample factor                 & \cv{rom.lstm.subsample}         & \cv{lstm\_rom.py}              \\
Residual prediction mode              & \cv{rom.lstm.residual}          & \cv{lstm\_rom.py}              \\
Patience (early stopping)             & \cv{rom.lstm.patience}          & \cv{lstm\_rom.py}              \\
Library polynomial degree $p_{\max}$  & \cv{wsindy.max\_poly}           & \cv{library.py}                \\
Test-function stride                  & \cv{wsindy.stride}              & \cv{system.py}                 \\
Morse forces enabled                  & \cv{forces.enabled}             & \cv{multifield.py}             \\
\bottomrule
\end{tabular}
\end{table}
